\section{Predicate operator syntax}

\begin{proposition}[Predicate syntax types]\label{thm:predicate-syntax-types}\lean{Obsmetrologia.PredicationSyntax.Ty}\leanok\uses{thm:entity-predicate}
The predication stage generates a typed operator grammar over the predicate
base: predicate type, function type, and product type. This is the usable
dependency-free core of the archived semiotic-universe syntax attempt.
\end{proposition}

\begin{proposition}[Intrinsically typed predicate terms]\label{thm:predicate-syntax-terms}\lean{Obsmetrologia.PredicationSyntax.Term}\leanok\uses{thm:predicate-syntax-types}
Variables and terms are indexed by their context and result type, so ill-typed
predicate operators are not representable in the active syntax carrier.
\end{proposition}

\begin{proposition}[Raw predicate operators]\label{thm:raw-predicate-operators}\lean{Obsmetrologia.PredicationSyntax.RawOp}\leanok\uses{thm:predicate-syntax-terms}
A raw $n$-ary predicate operator is a closed term from the right-associated
$n$-input predicate carrier into the predicate base.
\end{proposition}

\begin{theorem}[Unary raw operators expose the paired carrier]\label{thm:rawop-one-carrier}\lean{Obsmetrologia.PredicationSyntax.rawOp_one}\leanok\uses{thm:raw-predicate-operators}
The unary raw-operator carrier unfolds to a closed term
$(P \times P)\to P$, matching the archive's finite-product convention while
remaining inside the dependency-free predication spine.
\end{theorem}

\begin{theorem}[Semantic interpretation and fusion closure]\label{thm:predicate-syntax-interpretation-frontier}\notready\uses{thm:raw-predicate-operators, thm:closure-operator}
The next step needed is not a free Heyting ambient. It is an interpretation of
predicate syntax into the active valuation/settling carriers that commutes
with the derived closure operations, followed by a fusion closure that names
exactly the already available admissible behaviors.
\end{theorem}
